\begin{agradecimentos}

\begin{center}
\textbf{Arthur Carneiro Trindade}
\end{center}

Primeiramente, agradeço à minha família, que esteve ao meu lado durante toda essa jornada.

À minha mãe, Eloíde, que desempenhou os papéis de mãe e pai para mim e para minhas irmãs, e, mesmo diante das dificuldades que a vida nos impôs, nunca deixou de me apoiar e de acreditar em mim, principalmente nos momentos em que o peso parecia grande demais para carregar.

Às minhas irmãs, Érica e Estela, que acompanharam de perto meu esforço, as noites em claro e as frustrações, agradeço por todas as conversas, pelo apoio e por terem compartilhado comigo tantos momentos.

E aos meus avós, Marta, Madalena e Francisco, que cuidaram de mim e que, infelizmente, partiram ao longo dessa jornada. Gostaria muito que estivessem aqui para compartilhar este momento comigo e, embora não possam presenciar a conclusão desta etapa, eu levo comigo tudo o que me ensinaram.

Agradeço também à minha segunda família, os amigos que tive a felicidade de encontrar ao longo desses anos. Posso afirmar, com toda a certeza, que não teria conseguido chegar até aqui sem vocês, que conseguiram transformar uma jornada que muitas vezes pareceu interminável e dolorosa em uma caminhada tão leve e agradável.

Em especial, ao meu amigo e irmão, Miguel, com quem tive o privilégio de compartilhar tantos momentos, bons e ruins, ao longo desses anos, e com quem também tive a honra de desenvolver este trabalho. Sou profundamente grato por sua amizade, sua confiança e por tudo o que construímos juntos. Tenho certeza de que esta é apenas uma das muitas conquistas que ainda compartilharemos. E também ao meu amigo Bruno, que esteve presente nessa caminhada e com quem tive o prazer de compartilhar tantos momentos ao longo desses anos, pela amizade e por sua contribuição para transformar este projeto em realidade.

Agradeço à Universidade de Brasília, por todos esses anos de aprendizado. Apesar dos sacrifícios que fizeram parte dessa trajetória, estudar em uma instituição de tamanho renome como a UnB foi uma experiência que me transformou tanto pessoal quanto profissionalmente. Agradeço a todos os professores que fizeram parte da minha formação e que contribuíram para que eu chegasse até aqui.

Ao professor John Lenon Cardoso Gardenghi, que esteve presente em diferentes momentos da minha formação, como professor, orientador de estágio e, por fim, orientador deste Trabalho de Conclusão de Curso, agradeço pela confiança e por ter aceitado, de boa vontade, nos acompanhar nesta última etapa. À professora Milene Serrano e ao professor Maurício Serrano, que também fizeram parte da minha formação, agradeço por aceitarem participar da banca examinadora deste trabalho.

Por fim, agradeço aos amigos e colegas que conheci ao longo da minha trajetória profissional, nos diferentes lugares em que tive a oportunidade de trabalhar. Mesmo que nossas trajetórias tenham se cruzado por pouco tempo, cada experiência contribuiu para o meu desenvolvimento, não apenas como profissional, mas também como pessoa. Agradeço, ainda, a todos aqueles que, direta ou indiretamente, contribuíram para a realização deste projeto.

Encerro esta etapa da vida com orgulho da trajetória que percorri, com a certeza de que não sou mais o mesmo homem que ingressou na universidade há oito anos e de que não teria chegado até aqui sozinho.

A todos vocês, meu muito obrigado.

\newpage

\begin{center}
\textbf{Miguel Moreira da Silva de Oliveira}
\end{center}

Agradeço primeiramente à minha mãe, Maria Nilde. Obrigado, mãe, por sempre me dar as melhores condições de vida, sem nunca deixar faltar nada. Obrigado por me educar da melhor forma possível, para que eu buscasse meus sonhos e objetivos sempre com respeito e educação, e por me dar acesso às melhores escolas, o que foi fundamental para eu chegar onde estou hoje. Obrigado por todo o amor, o carinho, as broncas, os ensinamentos e por tudo o que você fez e faz por mim. Você é minha maior inspiração e eu admiro muito a sua história. Se não fosse pela sua dedicação à minha criação, eu não estaria aqui. Tudo isso é graças a você, e tudo o que faço é para te dar um futuro ainda melhor do que o passado e o presente que você me deu.

Ao meu irmão, Rafael Moreira, obrigado pela educação e pelos ensinamentos que você me dá desde sempre, e por me colocar de volta nos trilhos quando começo a me dispersar. Obrigado por todas as conversas: mesmo que você não entenda muito do que eu falo sobre software, me ouvir me ajuda a pensar. Obrigado pelas piadas e brincadeiras que, nos dias mais cheios e exaustivos, me ajudam a relaxar. Obrigado por ser o melhor irmão que eu poderia ter. E, no BoraFut, obrigado por sempre acreditar no projeto, me incentivar, sugerir, perguntar, ouvir e discutir para melhorar, e por ter feito parte da idealização de tudo lá no início, em 2023.

Ao Arthur Trindade, minha dupla neste TCC e, na verdade, minha dupla para a vida, em qualquer coisa e principalmente no software. Quando nos encontramos em MDS, quem diria que estaríamos aqui hoje entregando um TCC juntos, e com ele um projeto gigante. Mas projetos gigantes já viraram rotina para nós: ao longo da graduação participamos de diversos projetos juntos, sempre entregando tudo com maestria. E este não será o último; ainda temos muitas ideias mapeadas, cada uma no seu tempo. Obrigado por me acompanhar nessa trajetória e por sempre topar minhas ideias malucas. Você é, sem dúvida, o desenvolvedor que eu mais admiro. Muita coisa deu ``errado'' para a gente se juntar pela primeira vez, mas desde então tudo deu certo. É mais um projeto entregue, mais um capítulo encerrado juntos, e que venham muitos outros nesse grande livro da vida. Você vai ser para sempre minha dupla.

Ao professor John Lenon Cardoso Gardenghi, nosso orientador neste trabalho, mas não apenas isso. O John esteve presente na nossa graduação desde o início: foi nosso primeiro contato como orientador de estágio, depois nosso professor em Fundamentos de Arquitetura de Computadores e em Estruturas de Dados 2 e, por fim, nosso orientador de TCC. Ao longo de todo esse caminho, foi nosso mentor dentro da UnB. Sempre que precisávamos de ajuda, recorríamos a ele, e ele sempre fez de tudo para nos atender. Foi uma pessoa essencial na nossa graduação e, mais ainda, neste trabalho. Obrigado, professor.

Ao Ruy Barbosa, meu primeiro supervisor de estágio, que me deu a oportunidade de estagiar no Instituto Nacional de Tecnologia da Informação (ITI), uma das maiores viradas de chave da minha vida profissional. Quando ele me contratou, eu não sabia nada; ele apostou na minha vontade e no meu empenho em aprender, algo que deveria ser comum, mas que hoje é raro em um mercado que valoriza mais a experiência prévia do que a disposição para aprender. O Ruy me ensinou toda a base, não apenas técnica, mas também sobre como lidar com problemas, buscar soluções, me portar profissionalmente e entender como o mercado funciona. Sem isso, eu não teria chegado onde estou hoje. Obrigado, Ruy.

À Carla Cangussu, ao Guilherme Peixoto e ao Henrique Hida, meus primeiros companheiros de estágio, que por coincidência também estudaram aqui na FCTE. Obrigado por todo o carinho, cuidado e paciência comigo. Como meus veteranos, vocês me ensinaram tudo o que sabiam, não só sobre o trabalho, mas também sobre a faculdade, e hoje somos grandes amigos. Obrigado por tanto.

Ao Arthur Trindade, Bruno Ricardo, Diogo Bites, Henrique Hida, João Gabriel, Luis Eduardo e Rafael Moreira, a formação inicial do BoraFut, lá em 2023. Meus grandes amigos, vocês toparam entrar nessa comigo acreditando no meu sonho e foram de importância sem igual para este trabalho e para o projeto. Por falta de tempo, não conseguimos seguir naquela época, mas hoje, com o BoraFut em produção, devo agradecer a vocês, que acreditavam desde o começo e que desde então sempre me ajudaram, alguns de forma indireta e outros de forma bem direta. Obrigado pela amizade e por acreditarem que podemos melhorar o futebol em Brasília e no Brasil. O futuro é brilhante, meus amigos, e estaremos todos juntos.

Ao Gabriel Leite, meu amigo e irmão. Em tantas manhãs em que acordei cedo e tantas noites em que passei programando, você sempre me fez companhia e compartilhou comigo todos os sentimentos, bons ou ruins. Além disso, me ajudou a resolver diversas questões do projeto, como divulgação, estrutura e logística, e sempre me acompanha e topa tudo comigo. Obrigado.

A todos os meus amigos, com quem sei que posso contar sempre. Obrigado por me escutarem, me apoiarem, acreditarem nos meus sonhos e projetos e estarem presentes para celebrar cada conquista comigo. Obrigado também por me ensinarem, me corrigirem e me avisarem quando estou errado, e por estarem comigo nos melhores e nos piores momentos. Todos vocês participaram disto, mesmo que indiretamente e mesmo que não percebam: aquela conversa, aquela jogatina, aquele fut, tudo isso ajuda a recarregar as baterias e a tirar forças de onde não tem para fazer acontecer. Obrigado a todos.

A todos os professores que fizeram parte da minha formação, por cada aula, conselho e conhecimento compartilhado ao longo dessa caminhada.

Por fim, à Universidade de Brasília e à Faculdade de Ciências e Tecnologias em Engenharia (FCTE), pela formação de qualidade e por todas as oportunidades que tornaram esta conquista possível.

\end{agradecimentos}
